\begin{enumerate}

\item Obtener área, centroide, I\_z, I\_y e I\_yz sumando aportaciones de pared y booms.

\item Determinar los gradientes de tensión axial con ambos cortantes y el producto de inercia; integrar dq/ds en cada pared.

\item Imponer continuidad y equilibrio en los nudos; en contornos abiertos q=0 en extremos libres y en celdas cerradas imponer compatibilidad de giros.

\item Localizar el centro de cortadura anulando el torsor debido al flujo; calcular la tensión máxima por |q(s)|/t en paredes y puntos críticos.

\item Las leyes completas q(s) por pared se indican a continuación; la coordenada s se mide desde el nudo inicial de cada pared y el sentido firmado sigue las flechas del diagrama. Los resultados geométricos se obtienen con la misma idealización de paredes y larguerillos del enunciado.

\item Pared inferior\_1: q(s) = $\frac{52 s}{73} - \frac{907400}{29273}$ N/mm; s en mm desde su nudo inicial, con signo según el recorrido.

\item Pared inferior\_2: q(s) = $\frac{52 s}{73} - \frac{617890}{29273}$ N/mm; s en mm desde su nudo inicial, con signo según el recorrido.

\item Pared derecha: q(s) = $\frac{- 13 s^{2} + 780 s}{1095} + \frac{841750}{29273}$ N/mm; s en mm desde su nudo inicial, con signo según el recorrido.

\item Pared superior\_2: q(s) = $\frac{841750}{29273} - \frac{52 s}{73}$ N/mm; s en mm desde su nudo inicial, con signo según el recorrido.

\item Pared superior\_1: q(s) = $\frac{760760}{29273} - \frac{52 s}{73}$ N/mm; s en mm desde su nudo inicial, con signo según el recorrido.

\item Pared izquierda: q(s) = $\frac{13 s^{2} - 780 s}{1095} - \frac{907400}{29273}$ N/mm; s en mm desde su nudo inicial, con signo según el recorrido.

\item Pared comun: q(s) = $\frac{- 13 s^{2} + 780 s}{876} + \frac{1378650}{29273}$ N/mm; s en mm desde su nudo inicial, con signo según el recorrido.

\item Pared crítica: comun; s = 30 mm desde su nudo inicial.

\end{enumerate}

\noindent\textbf{Operaciones y resultados intermedios.}

\begin{enumerate}

\item Área equivalente que soporta tensión axial.

\[\fitmath{y_{1}=2\cdot 80+2\cdot 70+2\cdot 60+2\cdot 70+2\cdot 80+2\cdot 60+2.5\cdot 60\approx 990\,\mathrm{mm^{2}}}\]

\item Propiedades geométricas de la sección (centroide y momentos de inercia).

\[\fitmath{z_G=75.7575758\,\mathrm{mm},\quad y_G=30\,\mathrm{mm}}\]

\item Momentos de inercia del modelo en los ejes centroidales.

\[\fitmath{I_z=657000\,\mathrm{mm}^{4},\quad I_y=2478181.82\,\mathrm{mm}^{4},\quad I_{yz}=0\,\mathrm{mm}^{4}}\]

\item Centro de cortadura a partir del equilibrio de los flujos; coordenadas del modelo.

\[\fitmath{z_C=76.3775493\,\mathrm{mm},\quad y_C=30\,\mathrm{mm}}\]

\item Tensión máxima |q|/t en pared comun.

\[\fitmath{y_{5}=\frac{\left|\frac{1769625}{29273}\right|}{\frac{5}{2}}\approx 24.18099\,\mathrm{MPa}}\]

\end{enumerate}
